%%
% BIThesis 研究生学位论文模板 The BIThesis Template for Graduate Thesis
% This file has no copyright assigned and is placed in the Public Domain.
%%

\chapter{绪论}
\label{chap:intro}

\section{研究背景及意义}

近年来，以Transformer架构为基础的大语言模型（Large Language Models，LLMs）\cite{vaswani2017attention}在自然语言理解、逻辑推理、文本生成和代码生成等任务中取得了重要进展。以GPT-4系列\cite{ref2}、LLaMA系列\cite{ref3,ref4,ref5}和Qwen系列\cite{ref6,ref7,ref8}为代表的基础模型，已经能够在开放问答、数学推理、程序合成、信息抽取和复杂指令遵循等场景中表现出较强的生成式推理能力。上述模型为复杂任务自动化提供了统一的语言接口和推理基础，也为多个智能体围绕同一任务进行分工、协作和验证提供了可能。

然而，大语言模型本身并不能完全满足复杂现实任务的需求。首先，模型参数中固化的知识具有时间滞后性，难以及时覆盖不断变化的外部世界知识\cite{ref28}。其次，模型的上下文窗口即使不断扩展，仍然会受到输入长度、推理延迟和调用成本的限制；当上下文过长时，模型还容易出现“迷失在中间（Lost in the Middle）”现象，即对位于长文本中部的关键信息关注不足\cite{ref29,ref30}。再次，单次模型调用缺少与外部环境的持续交互能力，难以主动检索信息、执行代码、检查中间结果或根据反馈修正推理路径。上述限制使得单一大语言模型在软件工程、科学计算、开放环境问答和具身智能等任务中仍然存在可靠性不足、知识更新困难和错误难以自我修复等问题\cite{ref9,ref10,ref11}。

为了弥补单次模型调用在规划、反馈和外部执行方面的不足，研究者开始将大语言模型组织为智能体（Agent）\cite{ref12}。智能体通常以大语言模型为核心推理器，并在其外部加入任务规划、工具调用、记忆管理和反馈修正等机制。思维链（Chain of Thought，CoT）\cite{ref13}、自一致性（Self-Consistency）\cite{ref14}、思维树（Tree of Thoughts，ToT）\cite{ref15}和思维图（Graph of Thoughts，GoT）\cite{ref16}等方法，主要通过显式化或扩展推理路径提升复杂问题求解能力；Toolformer\cite{ref17}、ToolLLM\cite{ref18}和ReAct\cite{ref23}等工作进一步将推理与行动结合，使模型能够在“思考--行动--观察”的闭环中调用搜索、计算器、代码执行器和外部API等工具；Reflexion\cite{ref25}、Generative Agents\cite{ref26}和Voyager\cite{ref27}等研究则通过自然语言反馈、长时记忆和环境交互提升智能体在多轮任务中的持续改进能力。

智能体使大语言模型能够以更接近任务执行流程的方式持续感知、规划、执行和反馈。对于需要外部知识、复杂推理或长周期交互的任务，智能体可以把模型能力与外部工具能力结合起来，从而缓解单模型知识封闭、执行能力不足和缺少反馈的问题。但单智能体仍受限于单一角色视角、单条推理路径和有限的自我纠错能力。当任务需要多个专业角色协作、不同推理路径相互验证或多个工具结果共同汇总时，仅依靠单个智能体往往难以稳定获得可靠结果。

在此背景下，多智能体系统（Multi-Agent Systems，MAS）进一步引入角色分工、协同推理和信息共享机制\cite{ref31,ref32}。当任务同时涉及规划、检索、执行、验证和汇总等环节时，多个具备不同职责或工具权限的智能体可以围绕同一目标形成互补视角，并通过消息传递完成任务分解、候选答案生成、交叉检查和最终决策。与单智能体相比，多智能体系统的潜在优势并不只是增加模型调用次数，而是通过角色差异和协作式纠错降低单一路径推理失败的风险。

已有研究表明，多智能体系统在软件开发、开放式问题求解和复杂推理任务中具有应用价值\cite{ref9,ref35,ref36}。多智能体辩论等研究也说明，不同模型或不同角色之间的交叉质询能够改善事实性和推理鲁棒性\cite{ref31,ref34}。这些工作表明，角色分工、通信路径和信息汇总方式会影响系统性能。若通信机制设计合理，不同智能体能够补充彼此的知识、发现候选答案中的错误并提升最终输出的可靠性；若通信机制设计不充分，系统可能只是增加对话规模，额外上下文和重复讨论反而会抵消协作带来的收益。

多智能体协作的关键并非简单增加智能体数量，而是在合适的时机交换有效信息。从系统结构看，多智能体系统可以表示为一个带有节点和边的通信图：节点代表智能体，边代表信息传递路径。每个智能体在当前轮次生成的中间推理、候选答案或工具观察，需要通过通信边传递给后继智能体；上一轮协作产生的历史状态，也可能在下一轮继续作为上下文使用。因此，通信拓扑决定哪些智能体能够读取哪些上游输出，消息内容决定下游智能体实际接收到什么信息，跨轮记忆则决定历史推理如何影响后续协作。

不同通信结构决定了信息覆盖范围、协作顺序和成本上界。通信过少时，系统难以充分利用多角色互补；通信过多时，系统又会被冗余上下文和高昂调用成本拖累。因此，多智能体系统性能不仅取决于单个智能体能力，还高度依赖通信结构和通信内容是否高效。

现有多智能体系统在通信机制上仍以粗粒度设计为主。一方面，许多系统采用静态或人工预设拓扑，通信边在任务执行前即被固定，难以根据任务类型、推理阶段和工具使用情况动态调整。另一方面，智能体之间传递的消息往往是完整自然语言输出，包括中间推理、背景解释、候选答案、工具观察和格式化说明等内容。对于下游智能体而言，这些完整消息并不总是全部有用：某些内容与接收方角色无关，某些观点已经被多个上游智能体重复表达，某些历史信息在后续轮次中的贡献有限。

如图\ref{fig:communication_bottleneck}所示，当智能体数量和交互轮次增加时，多智能体系统中的输入词元（Token）数量和整体调用成本会随之增长。在完全连接拓扑下，若每一轮每个智能体都向其他智能体传递完整输出，则通信边数量会随智能体数量呈平方级增长；若系统还保留跨轮历史，则上下文长度会进一步随轮次累积。对于按Token计费的闭源模型和商业API而言，这种增长会导致部署成本和推理延迟上升；对于本地部署模型而言，过长上下文也会增加显存占用和推理时间。

\begin{figure}[H]
  \centering
  \makebox[\textwidth][c]{%
    \subcaptionbox{%
     多智能体系统的Token消耗
      \label{fig:communication_bottleneck_left}%
    }[0.45\textwidth]{%
      \includegraphics[width=0.45\textwidth]{\detokenize{image/第一章/图 1‑1-a.png}}%
    }%
    \hspace{0.02\textwidth}%
    \subcaptionbox{%
       不同智能体数量下的输入Token
      \label{fig:communication_bottleneck_right}%
    }[0.45\textwidth]{%
      \includegraphics[width=0.45\textwidth]{\detokenize{image/第一章/图 1‑1-b.png}}%
    }%
  }
  \caption{多智能体系统的Token消耗对比和不同智能体数量下的输入Token变化}
  \label{fig:communication_bottleneck}
\end{figure}

通信开销的增加并不一定带来任务性能提升。相反，冗长、重复和弱相关的上下文会降低下游智能体提取关键信息的效率。当模型需要在较长历史对话中寻找有效证据时，关键线索可能被无关解释、重复推理或过长工具返回结果淹没，从而引发长上下文注意力衰减问题\cite{ref29}。因此，多智能体通信优化需要同时回答两个问题：一条通信边上传递的内容是否有效，以及这条通信边本身是否仍然值得保留。

本文将上述通信瓶颈分解为两个相互关联但层次不同的问题。语义冗余问题发生在既定通信边上的消息内容内部，表现为同一条自然语言消息包含较多与接收方角色无关的解释、重复推理、弱相关背景或过长历史摘要。该问题关注在既定通信边上应该传递什么内容，哪些信息应被保留，哪些冗余文本可以被压缩。通信冗余问题发生在通信图结构层面，表现为某些智能体之间的连接本身贡献有限，甚至会传播高成本、低收益的工具观察和历史状态。该问题关注哪些智能体之间需要保持通信，哪些边应当被削弱或剪除。

这两个问题在工具调用场景中更加突出。具备工具使用能力的智能体不仅生成自然语言推理，还会调用搜索、代码执行、文本检查和外部API等工具\cite{ref17,ref18,ref23}。工具返回内容往往较长，并可能包含网页片段、运行日志、候选代码、异常信息或中间计算结果。在单智能体系统中，这些观察主要增加当前智能体的上下文；而在多智能体系统中，工具观察还可能沿通信边传播给多个下游智能体，触发重复阅读、重复分析甚至二次工具调用。因此，一条通信边的成本不再只由消息长度决定，还与发送方工具调用次数、工具返回Token和下游传播范围有关。

基于上述分析，本文将多智能体高效通信研究划分为内容层面和结构层面两条互补路线：内容层面通过消息压缩缓解语义冗余问题，提高既定通信边上传递信息的密度；结构层面通过图剪枝缓解通信冗余问题，减少低贡献通信边和高成本工具观察传播。二者共同服务于同一目标，即在保持或提升多智能体协作收益的同时，降低Token消耗、工具调用开销和整体推理成本。

研究多智能体系统中的高效通信机制具有明确的理论意义和实际价值。本文将多智能体系统中的通信效率问题明确拆分为“消息内容是否冗余”和“通信边是否必要”两个层次，为理解大语言模型多智能体协作中的成本--质量权衡提供了更清晰的分析框架。现有研究往往分别关注提示词压缩、长上下文处理或通信拓扑优化，但较少将自然语言消息内部冗余与通信图结构冗余放在统一视角下讨论。本文围绕语义冗余问题和通信冗余问题展开研究，有助于揭示多智能体系统中上下文膨胀、信息重复、工具观察传播和通信边贡献差异之间的关系。

在方法层面，本文旨在不改变底层大语言模型参数的前提下提升多智能体通信效率。对于闭源商业模型或难以微调的大模型系统而言，外部通信机制优化具有较强的通用性和可迁移性。消息压缩方法可以嵌入智能体之间的信息传递过程，使下游智能体接收信息密度更高的上下文；图剪枝方法则可以从通信结构角度选择贡献更高的边，使系统不必无差别地保留所有潜在连接。两类方法分别对应“传什么”和“向谁传”两个关键决策。

在实际应用层面，高效通信是多智能体系统实现大规模部署的重要前提。在工业和企业应用中，任务通常需要多轮推理、多工具调用和多角色协作，如果每个智能体都接收未筛选的完整历史消息，系统成本会随任务复杂度上升。通过减少冗余Token数量、压缩低贡献历史信息、削弱高成本低收益通信边，本文研究能够降低API调用开销和推理延迟，提高多智能体系统在实际应用中的经济可行性。同时，经过筛选和压缩的通信内容能够减少无关信息干扰，使下游智能体更便于定位关键证据和主要分歧，从而有助于提升任务完成质量和系统稳定性。

\section{国内外研究现状}

随着大语言模型从单次问答扩展到多轮任务执行，国内外研究逐渐围绕智能体协作、上下文压缩和通信结构优化形成了较为清晰的研究脉络。总体来看，相关工作可以分为三条路线：第一类研究关注如何将多个大语言模型调用组织为可协作的智能体系统，重点解决复杂任务中的角色分工、对话调度和工具使用问题；第二类研究关注如何压缩冗长提示词或对话历史，重点减少输入上下文中的低贡献语义内容；第三类研究关注如何优化多智能体通信图或工作流结构，重点减少低贡献智能体、冗余通信边和无效协作路径。本节从上述三条路线概括国内外研究现状，具体方法分类将在第\ref{chap:related-work}章进一步展开。

在多智能体协作框架方面，已有研究首先证明了将多个智能体组织起来可以提升复杂任务处理能力。多智能体辩论方法通过让多个模型或多个角色给出独立推理，再经由辩论、投票或评审机制形成最终答案，从而提升事实性和推理鲁棒性\cite{ref31,ref34}。面向软件工程和复杂任务执行的ChatDev\cite{ref9}、MetaGPT\cite{ref35}和AutoGen\cite{ref36}等框架，则进一步将角色分工、流程控制、工具调用和结果汇总纳入统一协作过程。CAMEL\cite{ref37}、AgentVerse\cite{ref79}和ChatEval\cite{ref80}等工作从角色扮演、智能体社会和多智能体评估角度探索多个智能体之间的交互方式。这些研究推动了多智能体系统从简单对话走向结构化任务执行，也说明系统性能不仅取决于单个模型能力，还与角色设置、通信方式和信息汇总策略密切相关。

然而，现有多智能体框架为了保证协作信息充分共享，往往倾向于保留较长历史上下文，或采用较密集的通信拓扑。对于小规模任务，这种设计能够降低系统实现难度；但当智能体数量、交互轮次和工具调用次数增加后，完整消息广播会增加Token成本和推理延迟。同时，冗长历史并不总是带来更多有效信息，弱相关解释、重复推理和过期状态还可能干扰下游智能体定位关键证据。因此，在多智能体框架逐渐成熟之后，如何提高通信效率成为进一步提升系统可用性的关键问题。

在语义内容压缩方面，国内外研究主要围绕长上下文压缩和提示词压缩展开。SelectiveContext\cite{ref55}通过词元或短语的信息量估计过滤低信息量内容，LLMLingua\cite{ref56}进一步提出从粗到细的提示词压缩框架，LLMLingua-2\cite{ref57}将压缩过程转化为词元级保留预测问题，以提升压缩效率和跨任务泛化能力。后续研究进一步从长上下文场景、软提示压缩和任务反馈等角度提升压缩质量\cite{ref85,ref86,ref58}。此外，检索增强与选择性压缩方法也尝试在外部证据进入模型前筛选信息，以减少冗余上下文带来的推理负担\cite{ref88,ref91,ref92}。这类方法为降低长上下文推理成本提供了重要基础，也为多智能体消息压缩提供了可借鉴的技术路径。

但将通用提示词压缩方法应用于多智能体通信仍存在不足。普通长文本压缩通常把输入视为一个整体文本，主要判断哪些词元、短语或句子应被保留；而多智能体通信消息具有明确的发送方、接收方、前驱节点和轮次状态。同一段推理对求解智能体可能有用，对验证智能体可能只需保留结论和证据；同一事实若已由多个前驱节点重复表达，继续保留全部表述只会增加上下文长度；跨轮历史中真正影响后续推理的也往往不是完整对话，而是共识、分歧和关键依据。因此，多智能体场景中的语义压缩需要结合角色关系、节点关系和历史状态，而不能仅以文本长度缩减作为目标。

在通信结构优化方面，相关研究从图结构学习、图剪枝和工作流搜索等角度减少无效通信。早期深度多智能体通信研究已经从可学习消息、注意力通信和何时通信等角度探索协作结构\cite{ref93,ref94,ref95}。在大语言模型多智能体系统中，DyLAN\cite{ref39}通过动态选择参与求解的智能体减少无效讨论轮次，CommFormer\cite{ref40}、G-Designer\cite{ref41}和MCGD\cite{ref42}等工作从图建模或图生成角度学习通信拓扑。AgentPrune\cite{ref43}将多智能体系统建模为空间--时间通信图，并通过可学习掩码剪除低贡献通信边；AgentDropout\cite{ref44}、SafeSieve\cite{ref45}和AGP\cite{ref46}则从节点丢弃、渐进式剪枝和自适应稀疏化等角度进一步优化通信结构。与此同时，GPTSwarm\cite{ref48}、AFlow\cite{ref50}和ADAS\cite{ref99}等方法把智能体系统视为可搜索或可优化的工作流，通过拓扑搜索或架构搜索获得更高效的协作路径。

上述通信结构优化方法能够有效减少部分低贡献节点和通信边，但在工具调用型多智能体系统中仍面临新的成本建模问题。具备工具使用能力的智能体不仅产生自然语言推理，还可能调用搜索、代码执行或文本检查工具；工具增强语言模型和函数调用框架已经表明，外部工具会改变模型的上下文组织方式和任务执行流程\cite{ref75,ref76,ref77}。这些工具返回内容常常包含网页片段、运行日志、异常信息或中间证据，长度和贡献差异较大。当这类工具观察沿通信边传播时，一条边的成本不再只由是否存在决定，还与发送方工具调用次数、工具返回Token和下游传播范围有关。因此，如果图剪枝方法只依据最终任务效用或边采样概率进行判断，就难以区分低成本高贡献边与高成本低收益边。

综合来看，现有研究已经分别从协作框架、语义内容压缩和通信结构优化三个方向推动了大语言模型多智能体系统的发展，但仍缺少对“消息内容冗余”和“通信边冗余”的统一分析。本文不以训练更强的基座模型或重构完整智能体框架为目标，而是在现有多智能体系统之上研究可插拔的通信优化方法：一方面，从消息内容层面压缩既定通信边上传递的自然语言信息，使下游智能体接收更短、更集中、更符合角色需求的上下文；另一方面，从通信图结构层面学习不同边在任务正确性和工具成本之间的权衡，减少低贡献通信路径和高成本工具观察传播。

因此，本文将普通多智能体推理任务和工具调用增强任务分别作为两个实验场景。普通推理任务用于检验消息级语义压缩是否能够提高上下文信息密度，工具调用增强任务用于检验通信边成本异质性下的图剪枝效果。本文在评价中同时关注任务准确率、Token成本、工具调用次数和工具Token，避免将通信优化简单等同于减少文本长度或减少边数量。只有在基本保持任务质量的前提下降低成本，或在相近成本下获得更高准确率时，多智能体通信优化才具有实际意义。

\section{本文主要研究内容及贡献}

本文围绕多智能体系统中的高效通信技术展开研究，目标是在不改动底层大语言模型参数的前提下，通过消息级压缩和图结构剪枝提升多智能体系统的准确率--成本平衡。总体思路是：首先分析多智能体系统中通信成本和上下文膨胀的来源，将其归纳为语义冗余问题和通信冗余问题；随后分别从通信内容和通信结构两个层面提出方法；最后在多类知识理解、数学推理、代码生成和工具调用任务上验证方法有效性。

为明确本文研究边界，本文将多智能体通信中的基础对象定义如下。设多智能体系统由若干智能体节点构成，通信拓扑表示智能体之间的信息传递路径。若同一轮中某个智能体输出被传递给另一个智能体，则该路径构成空间通信边；若上一轮智能体状态被保留并传递到下一轮，则构成时间通信边。智能体在边上传递的内容称为语义消息，通常包含中间推理、候选答案、工具观察和历史摘要等信息。

\begin{figure}[H]
  \centering
  \includegraphics[width=0.50\textwidth]{\detokenize{image/第一章/图 1‑2：完全通信拓扑.png}}
  \caption{完全连接拓扑（Full Topology）}
  \label{fig:full_Topology}
\end{figure}

在完全拓扑中，每个智能体都可以接收其他智能体的输出，如图\ref{fig:full_Topology}所示。这种拓扑能够保证信息覆盖，但也最容易产生冗余通信。结合前文分析，本文聚焦以下两个研究问题。

\textbf{问题1.1语义冗余问题：}在多智能体持续协作过程中，随着交互轮次增加和上游节点增多，智能体之间传递的自然语言消息会出现角色不匹配、观点重复和历史累积等问题。如何根据接收方角色、前驱节点差异和跨轮历史状态，筛选并保留对后续推理有效的信息，从而在降低Token成本的同时保持任务准确率？

\textbf{问题1.2通信冗余问题：}在工具调用型多智能体系统中，不同通信边携带的模型输出、工具调用次数和工具观察Token差异较大。如何将通信边贡献建模为可学习概率，并结合任务正确性与工具调用成本判断低贡献通信边，从而动态压缩通信图结构？

\subsection{面向语义冗余的消息压缩框架}

针对既定通信边上自然语言消息冗长、重复和弱相关的问题，本文提出面向语义冗余的消息压缩框架。该工作关注“通信时应传递什么内容”，目标是在不改变原有多智能体拓扑和底层模型的前提下，对智能体之间传递的推理内容进行细粒度筛选与聚合。

具体而言，本文设计了角色感知压缩、节点感知聚合和分歧感知压缩三个模块。角色感知压缩根据接收方智能体职责过滤发送方消息中的无关推理内容，解决完整消息与接收方需求不匹配的问题；节点感知聚合对多个前驱智能体传来的相似观点进行去重和来源标注，解决多前驱消息重复的问题；分歧感知压缩从跨轮历史中提取共识答案、候选分歧及其理由，解决多轮协作中历史上下文持续膨胀的问题。通过上述设计，本文方法能够在保留关键语义信息的同时减少冗余Token数量，提高下游智能体接收上下文的信息密度。

在实验方面，本文在MMLU、GSM8K、AQuA、MultiArith、SVAMP和HumanEval等任务上评估消息压缩框架，并与Direct、FullConnected和AgentPrune等基线进行对比。实验结果表明，该方法能够在多数任务上保持或提升任务准确率，同时相对于完全通信拓扑降低推理成本。

\subsection{面向通信边冗余的图剪枝算法}

针对工具调用场景中通信边成本差异较大的问题，本文提出面向通信边冗余的图剪枝算法。该工作关注“哪些智能体之间需要通信”，目标是在搜索、代码执行和文本检查等工具增强任务中，识别并剪除低贡献、高成本的通信边，从图结构层面降低多智能体系统的冗余通信开销。

具体而言，本文将智能体通信边建模为可学习概率，在训练过程中结合任务正确性、模型输出Token、工具调用次数和工具返回Token构造边级优化目标。与只根据最终任务效用评估通信结构的方法不同，本文进一步将工具调用带来的额外成本归因到具体发送方和通信边，使系统能够区分高贡献通信边与传播冗余工具观察的低贡献通信边。在此基础上，本文引入Taylor一阶重要性评估，根据边参数与梯度共同衡量移除该边对目标函数的影响，并对低重要性边进行动态剪枝。

在实验方面，本文在AIME、GAIA、GPQA、HLE、MATH500、MuSiQue和BigCodeBench等工具调用任务上验证面向通信边冗余的图剪枝算法的有效性。实验结果表明，该方法能够在工具调用场景中更合理地平衡任务准确率与推理成本，并在多个任务上相对完全通信或原始AgentPrune基线减少冗余工具信息传播。

\section{论文组织结构}

本文后续章节组织如下：第\ref{chap:related-work}章介绍多智能体系统、语义冗余问题和通信冗余问题相关研究，分析现有方法的优势与不足；第3章围绕面向语义冗余的消息压缩框架展开，给出系统执行流程、核心模块设计以及实验结果分析；第\ref{chap:tool-aware-pruning}章进一步面向工具调用场景，介绍面向通信边冗余的图剪枝算法及Taylor一阶重要性剪枝策略；最后，总结全文工作并讨论后续研究方向。
